\subsection*{2.2. 一种相容性（迹态射）}

\subsubsection*{2.2.1.}
设 $S$ 为 Hensel 离散赋值环谱, $u:X_s\longrightarrow Y_s$ 为 $s$-概形间有限表现的分离、拟有限、平坦态射. 对 $Y_s$ 上的每个阿贝尔层 $F$, 有迹态射
\[
\operatorname{Tr}_u:u_!u^*F\longrightarrow F
\]
（SGA~4, XVII, 6.2.3）. 若 $F$ 为 $Y_s\times_sS$ 或 $Y_s\times_s\eta$ 上的阿贝尔层, 则由此得到类似的态射 $\operatorname{Tr}_u:u_!u^*F\longrightarrow F$.

现在设 $u:X\longrightarrow Y$ 为 $S$-概形间有限表现的分离、拟有限、平坦态射. 仍用 $u$ 表示所诱导的态射
\[
X_s\times_sS\longrightarrow Y_s\times_sS,
\]
并对 $X$ 和 $Y$ 都用 $\Psi$ 表示 1.3.3 中的函子.

\subsubsection*{命题 2.2.2.}
\underline{对 $Y$ 上的每个阿贝尔层 $F$, 下图交换:}
\begin{center}
\begin{tikzcd}[column sep=large, row sep=large]
u_!u^*\Psi F
  \arrow[r, "{(1.3.7.1)}"] \arrow[d, "\operatorname{Tr}_u"']
  & u_!\Psi u^*F \arrow[r, "{(1.3.8.1)}"]
  & \Psi u_!u^*F \arrow[d, "\operatorname{Tr}_u"] \\
\Psi F \arrow[rr, equal] & {} & \Psi F.
\end{tikzcd}
\end{center}
验证留给读者. 对 $\Psi_\eta$ 及 $Y_\eta$ 上的 $F$ 也有相应变体.

\subsubsection*{2.2.3.}
设 $\Lambda$ 如 2.1.1, 并设 $u:X\longrightarrow Y$ 为分离有限型平坦 $S$-概形间的拟有限平坦态射:
\begin{center}
\begin{tikzcd}[column sep=large, row sep=large]
X \arrow[rr, "u"] \arrow[dr, "f"'] & {} & Y \arrow[dl, "g"] \\
{} & S & {}
\end{tikzcd}
\end{center}
原印本此处指向 2.2.1, 但系数环 $\Lambda$ 是在 2.1.1 中固定的. 仍以 $f$ 和 $g$ 表示态射
\[
X_s\times_sS\longrightarrow S,
\qquad
Y_s\times_sS\longrightarrow S,
\]
并假设 $\dim Y_\eta\leq n$. 对每个 $K\in\operatorname{Ob}D(S,\Lambda)$, 有迹态射（SGA~4, XVIII, 2.9）
\[
\operatorname{Tr}_f:Rf_!f^*K(n)[2n]\longrightarrow K.
\]

% Independent reference audit: printed/source (2.1.10.1) and (2.1.10.2)
% do not exist in this expose. Both arrows and the functoriality statement
% concern Rf_! R\Psi -> Rf_!, which is (2.1.8.2).
\subsubsection*{推论 2.2.4.}
\underline{下图交换:}
\begin{center}
\resizebox{\textwidth}{!}{%
\begin{tikzcd}[ampersand replacement=\&, column sep=large, row sep=large]
Rg_!u_!u^*R\Psi\,g^*K(n)[2n]
  \arrow[r, "{(1.3.7.1)}"] \arrow[d, "\operatorname{Tr}_u"']
  \& Rf_!R\Psi\,f^*K(n)[2n]
      \arrow[r, "{(2.1.8.2)}"]
  \& Rf_!f^*K(n)[2n] \arrow[d, "\operatorname{Tr}_f"] \\
Rg_!R\Psi\,g^*K(n)[2n]
  \arrow[r, "{(2.1.8.2)}"]
  \& Rg_!g^*K(n)[2n] \arrow[r, "\operatorname{Tr}_g"]
  \& K
\end{tikzcd}}
\end{center}
\underline{若 $u$ 平展, 则 (1.3.7.1) 为同构.}\EditorialFootnote{法文印本及英文工作版在上图将相关箭头标为 (2.1.10.1), 并在下文引用 (2.1.10.2); 本讲没有这两个编号. 此处的态射及其函子性对应 (2.1.8.2).}

省略自变量后, 图 (2.2.4) 可展开如下:
\begin{center}
\resizebox{\textwidth}{!}{%
\begin{tikzcd}[ampersand replacement=\&, column sep=large, row sep=huge]
Rg_!u_!u^*R\Psi\,g^*
  \arrow[r] \arrow[d, "\operatorname{Tr}_u"']
  \arrow[rd, phantom, "{(1)}" description]
  \& Rg_!R\Psi\,u_!u^*g^*
      \arrow[r] \arrow[d, "\operatorname{Tr}_u"']
      \arrow[rd, phantom, "{(2)}" description]
  \& Rg_!u_!u^*g^*
      \arrow[r, "\sim"] \arrow[d, "\operatorname{Tr}_u"']
      \arrow[rd, phantom, "{(3)}" description]
  \& Rf_!f^* \arrow[d, "\operatorname{Tr}_f"] \\
Rg_!R\Psi\,g^* \arrow[r, equal]
  \& Rg_!R\Psi\,g^* \arrow[r]
  \& Rg_!g^* \arrow[r, "\operatorname{Tr}_g"]
  \& K.
\end{tikzcd}}
\end{center}
(1) 的交换性由 2.2.2 推出; (2) 的交换性就是 (2.1.8.2) 的函子性; (3) 的交换性就是迹的传递性（SGA~4, XVIII, 2.9, 变体 3）.

\subsubsection*{2.2.4.1.}
对 $\Psi_\eta$ 和 $K\in\operatorname{Ob}D(\eta,\Lambda)$ 有类似结果.

\subsubsection*{2.2.5.}
设 $\Phi$ 为 $Y_s$ 的一个在 $s$ 上固有的子集, 以 $Rg_\Phi$ 表示“支集在 $\Phi$ 中的正像”函子. 以 $\operatorname{Tr}_g$ 表示复合
\begin{equation}\tag{2.2.5.1}
\begin{split}
\operatorname{Tr}_g:
Rg_\Phi R\Psi g^*K(n)[2n]
&\longrightarrow Rg_!R\Psi g^*K(n)[2n] \\
&\longrightarrow Rg_!g^*K(n)[2n]
\longrightarrow K,
\end{split}
\end{equation}
对 $\Psi_\eta$ 和 $K\in D(\eta,\Lambda)$ 也作类似定义. 原印本在 (2.2.5.1) 中将第一个函子写为 $Rf_\Phi$; 紧前的定义以及复合的其余部分要求这里为 $Rg_\Phi$, 故采用此写法.

若 $\Phi$ 为有限闭点集, 则
\[
Rg_\Phi\circ u_!=Rf_{u^{-1}\Phi},
\]
因而有迹态射
\[
\operatorname{Tr}_u:
Rf_{u^{-1}\Phi}u^*=Rg_\Phi u_!u^*\longrightarrow Rg_\Phi.
\]
由 2.2.4, 下图交换:
\begin{center}
\resizebox{\textwidth}{!}{%
\begin{tikzcd}[ampersand replacement=\&, column sep=large, row sep=large]
Rf_{u^{-1}\Phi}u^*R\Psi\,g^*K(n)[2n]
  \arrow[r] \arrow[d, "\operatorname{Tr}_u"']
  \& Rf_{u^{-1}\Phi}R\Psi\,f^*K(n)[2n]
      \arrow[r, "\operatorname{Tr}_f"]
  \& K \arrow[d, equal] \\
Rg_\Phi R\Psi\,g^*K(n)[2n]
  \arrow[rr, "\operatorname{Tr}_g"']
  \& {} \& K.
\end{tikzcd}}
\end{center}

假设 $u$ 平展, 取 $x\in X_s$, 并假设 $x$ 与 $y=u(x)$ 有相同的剩余域. 则
\[
Rf_{\{x\}}\xrightarrow{\ \sim\ }Rg_{\{y\}}.
\]
原印本在右侧也写成 $\{x\}$, 但 $g$ 的支集位于 $Y_s$ 中, 应为 $\{y\}$, 正如下图所示. 由 $\operatorname{Tr}_u$ 给出的这个同构表达了支集在 $x$ 处的上同调的平展局部性. 图
\begin{equation}\tag{2.2.5.2}
\begin{tikzcd}[column sep=huge, row sep=large]
Rf_{\{x\}}R\Psi\,f^*K(n)[2n]
  \arrow[r, "\operatorname{Tr}_f"] \arrow[d, equal]
  & K \arrow[d, equal] \\
Rg_{\{y\}}R\Psi\,g^*K(n)[2n]
  \arrow[r, "\operatorname{Tr}_g"]
  & K
\end{tikzcd}
\end{equation}
交换.

\subsubsection*{附注 2.2.6.}
相容性 (2.2.5.2) 表明, 当 $\Phi=\{y\}$ 时, 迹态射 (2.2.5.1) 是局部不变量: 它只依赖于 $S$-概形 $X$ 在 $x$ 处的 Hensel 化. 对 $R\Psi_\eta$ 也成立. 原印本对 (2.3.5.2) 与 (2.3.5.1) 的引用实指紧前的 (2.2.5.2) 与 (2.2.5.1); 同样, 由 $\Phi\subset Y_s$ 可见相关支集为 $\{y\}$, 而非原印本的 $\{x\}$.
